\documentclass[12pt, a4paper]{article}

% 1. Cấu hình Tiếng Việt và Font Times New Roman 13pt chuẩn
\usepackage[utf8]{inputenc}
\usepackage[T5]{fontenc}
\usepackage[vietnamese]{babel}
\usepackage{newtxtext,newtxmath}

\makeatletter
\renewcommand{\normalsize}{\@setfontsize\normalsize{13pt}{16pt}}
\makeatother
\normalsize

% 2. Gói Toán học & Ký hiệu bổ trợ
\usepackage{amsmath, amssymb, amsfonts, mathrsfs, amsthm}
\usepackage{graphicx}
\usepackage{fontawesome}
\usepackage{booktabs, tabularx, array, multirow}
\usepackage{enumitem}

% 3. Đồ họa TikZ, Bảng biến thiên & Biểu đồ
\usepackage{tikz}
\usetikzlibrary{calc, intersections, angles, quotes, patterns, arrows.meta, 3d}
\usepackage{pgfplots}
\pgfplotsset{compat=1.18}
\usepackage{tkz-euclide}
\usepackage{tkz-tab}

\renewcommand{\vec}[1]{\overrightarrow{#1}}

% 4. Hộp sư phạm (Tcolorbox)
\usepackage[many]{tcolorbox}
\tcbuselibrary{skins, breakable}

\newtcolorbox{pedabox}[1][]{
    enhanced,
    colback=blue!5!white,
    colframe=blue!75!black,
    fonttitle=\bfseries\sffamily,
    title=#1,
    boxrule=0.8pt,
    arc=2mm,
    breakable,
    left=3mm, right=3mm, top=2mm, bottom=2mm
}

\newtcolorbox{scaffoldbox}[1][]{
    enhanced,
    colback=orange!5!white,
    colframe=orange!85!black,
    fonttitle=\bfseries\sffamily,
    title=#1,
    boxrule=0.8pt,
    arc=2mm,
    breakable,
    left=3mm, right=3mm, top=2mm, bottom=2mm
}

\newtcolorbox{aibox}[1][]{
    enhanced,
    colback=green!5!white,
    colframe=teal!80!black,
    fonttitle=\bfseries\sffamily,
    title=#1,
    boxrule=0.8pt,
    arc=2mm,
    breakable,
    left=3mm, right=3mm, top=2mm, bottom=2mm
}

% 5. Căn lề chuẩn văn bản giáo án
\usepackage{geometry}
\geometry{
    a4paper,
    left=25mm,
    right=15mm,
    top=20mm,
    bottom=20mm
}

% 6. Header / Footer chuẩn hóa (BẮT BUỘC CHÍNH XÁC NỘI DUNG NÀY)
\usepackage{fancyhdr}
\usepackage{lastpage}
\pagestyle{fancy}
\fancyhf{}

\fancyhead[L]{\small \textit{Kế hoạch bài dạy Toán 11}}
\fancyhead[R]{\textbf{Trang \thepage/\pageref{LastPage}}}
\fancyfoot[L]{\small \textit{Tổ Toán - Tin}}
\fancyfoot[R]{\small \textit{Trường THCS\&THPT Hồng Vân}}
\renewcommand{\headrulewidth}{0.4pt}
\renewcommand{\footrulewidth}{0.4pt}

% 7. Giãn dòng & Giãn đoạn theo chuẩn Nghị định 30/2020/NĐ-CP
\usepackage{setspace}
\setstretch{1.15}
\setlength{\parindent}{1.0cm}
\setlength{\parskip}{5pt}

\begin{document}

\begin{center}
    {\Large \textbf{KẾ HOẠCH BÀI DẠY MÔN TOÁN LỚP 11}}\\[0.4em]
    {\LARGE \textbf{KIỂM TRA GIỮA HỌC KÌ I}}\\[0.5em]
    \textbf{Môn học:} Toán 11 | \textbf{Thời lượng:} 02 tiết (Tiết 23 và Tiết 24 theo PPCT)
\end{center}

\vspace{0.5em}
\hrule
\vspace{1em}

\section*{I. MỤC TIÊU ĐÁNH GIÁ}

\subsection*{1. Kiến thức, Năng lực}

\textbf{a) Năng lực Toán học đặc thù (nguyên văn chuẩn CT GDPT 2018):}
\begin{itemize}[leftmargin=1.5em]
    \item \textbf{Năng lực tư duy và lập luận toán học:}
    \begin{itemize}
        \item Đánh giá mức độ nhận biết và thông hiểu các kiến thức cơ bản từ đầu năm học đến hết Chương III:
        \begin{itemize}
            \item \textbf{Chương I (Hàm số lượng giác và phương trình lượng giác):} Nhận biết khái niệm góc lượng giác, giá trị lượng giác của một góc; mô tả các hệ thức lượng giác cơ bản, quan hệ góc liên quan đặc biệt; giải thích công thức biến đổi lượng giác; nhận diện đồ thị, tập xác định, tính chẵn lẻ, chu kì của 4 hàm lượng giác; nhận diện công thức nghiệm của phương trình lượng giác cơ bản.
            \item \textbf{Chương II (Dãy số. Cấp số cộng và cấp số nhân):} Nhận biết dãy số tăng, giảm, bị chặn; nhận biết và giải thích định nghĩa, tính chất, số hạng tổng quát, công thức tính tổng $n$ số hạng đầu tiên của cấp số cộng ($u_n = u_1 + (n - 1)d, S_n = \dfrac{n(u_1 + u_n)}{2}$) và cấp số nhân ($u_n = u_1 \cdot q^{n-1}, S_n = \dfrac{u_1(1 - q^n)}{1 - q}$).
            \item \textbf{Chương III (Các số đặc trưng đo xu thế trung tâm của mẫu số liệu ghép nhóm):} Hiểu ý nghĩa việc ghép nhóm số liệu; nhận biết và tính toán các số đặc trưng đo xu thế trung tâm: số trung bình cộng ($\overline{x}$), trung vị ($M_e$), tứ phân vị ($Q_1, Q_2, Q_3$), mốt ($M_o$) và khoảng tứ phân vị ($\Delta_Q = Q_3 - Q_1$).
        \end{itemize}
        \item Phân tích và lập luận logic trong các tình huống biến đổi toán học tổng hợp; suy luận loại trừ phương án nhiễu trong trắc nghiệm khách quan.
    \end{itemize}
    \item \textbf{Năng lực giải quyết vấn đề toán học:}
    \begin{itemize}
        \item Vận dụng các công thức lượng giác rút gọn biểu thức, giải phương trình lượng giác tìm nghiệm thuộc khoảng hoặc thỏa mãn điều kiện đại số.
        \item Thiết lập và giải hệ phương trình xác định các yếu tố của cấp số cộng, cấp số nhân; tính tổng các cấp số trong các bài toán đại số và thực tế.
        \item Đọc hiểu bảng tần số ghép nhóm, xác định nhóm chứa trung vị, tứ phân vị, mốt và tính toán chính xác giá trị bằng công thức nội suy.
        \item Quản lý và phân bổ thời gian làm bài thi 90 phút hiệu quả giữa phần trắc nghiệm (70\%) và tự luận (30\%).
    \end{itemize}
    \item \textbf{Năng lực mô hình hóa toán học:}
    \begin{itemize}
        \item Vận dụng kiến thức lượng giác, cấp số và thống kê để mô hình hóa và giải quyết bài toán thực tế: dao động điều hòa, lãi suất ngân hàng, thống kê số liệu y tế học đường.
    \end{itemize}
    \item \textbf{Năng lực giao tiếp toán học:}
    \begin{itemize}
        \item Trình bày bài làm tự luận chặt chẽ, mạch lạc; lập luận rõ ràng các bước biến đổi; sử dụng chuẩn xác các kí hiệu toán học $u_n, d, q, S_n, \overline{x}, M_e, M_o, Q_k$.
    \end{itemize}
    \item \textbf{Năng lực sử dụng công cụ, phương tiện học toán:}
    \begin{itemize}
        \item Sử dụng thành thạo máy tính cầm tay (MTCT Casio fx-580VN X) để kiểm tra kết quả tính toán, bấm máy giải phương trình, tính giá trị lượng giác và các số đặc trưng thống kê.
    \end{itemize}
\end{itemize}

\textbf{b) Năng lực số (theo Thông tư 02/2025/TT-BGDĐT):}
\begin{itemize}[leftmargin=1.5em]
    \item \textbf{Mã 2.6NC1a (Quản lý dữ liệu và an toàn thông tin kiểm tra):} Học sinh thực hiện đúng quy chế thi cử số hóa; nhận thức rõ yêu cầu bảo mật đề thi, mã đề và thông tin cá nhân trong quá trình chấm thi tự động và lưu trữ điểm số trên hệ thống cơ sở dữ liệu ngành giáo dục.
\end{itemize}

\textbf{c) Năng lực AI (theo Quyết định 2422/QĐ-BGDĐT):}
\begin{itemize}[leftmargin=1.5em]
    \item \textbf{Mã 11.B2.1 (Đạo đức học thuật và cam kết liêm chính trong kiểm tra đánh giá):} Học sinh thể hiện tính trung thực tuyệt đối trong phòng thi; không sử dụng bất kì thiết bị công nghệ hay ứng dụng AI nào để gian lận; hiểu rằng bài kiểm tra trực tiếp trên giấy là thước đo chuẩn xác năng lực tư duy toán học độc lập của chính mình.
\end{itemize}

\textbf{d) Phẩm chất:}
\begin{itemize}[leftmargin=1.5em]
    \item \textbf{Trung thực:} Tự giác làm bài, trung thực trong thi cử, không quay cóp hoặc trao đổi bài.
    \item \textbf{Chăm chỉ, trách nhiệm:} Tập trung cao độ trong suốt 90 phút làm bài, kiểm tra cẩn thận trước khi nộp bài.
\end{itemize}

\section*{II. HÌNH THỨC, MA TRẬN VÀ BẢN ĐẶC TẢ ĐỀ KIỂM TRA}

\subsection*{1. Hình thức kiểm tra}
\begin{itemize}[leftmargin=1.5em]
    \item \textbf{Hình thức:} Kết hợp giữa Trắc nghiệm khách quan (70\%) và Tự luận (30\%).
    \item \textbf{Thời gian làm bài:} 90 phút (2 tiết kiểm tra liên tiếp).
    \item \textbf{Cấu trúc điểm số:}
    \begin{itemize}
        \item Phần I: Trắc nghiệm nhiều lựa chọn gồm 35 câu (mỗi câu 0,2 điểm) = 7,0 điểm.
        \item Phần II: Tự luận gồm 4 bài toán = 3,0 điểm.
    \end{itemize}
\end{itemize}

\subsection*{2. Ma trận đề kiểm tra giữa học kì I}

\begin{center}
\begin{tabularx}{\linewidth}{|p{3.2cm}|c|c|c|c|c|}
\hline
\textbf{Chủ đề / Chương} & \textbf{Nhận biết} & \textbf{Thông hiểu} & \textbf{Vận dụng} & \textbf{Vận dụng cao} & \textbf{Tổng số} \\
\hline
\textbf{Chương I: Lượng giác} & 6 câu TN (1,2đ) & 4 câu TN (0,8đ) & 2 câu TN + 1 TL (1,2đ) & 1 câu TL (0,5đ) & 12 TN + 2 TL (3,7đ) \\
\hline
\textbf{Chương II: Cấp số} & 6 câu TN (1,2đ) & 4 câu TN (0,8đ) & 2 câu TN + 1 TL (1,1đ) & 1 câu TL (0,5đ) & 12 TN + 2 TL (3,6đ) \\
\hline
\textbf{Chương III: Thống kê} & 4 câu TN (0,8đ) & 5 câu TN (1,0đ) & 2 câu TN + 1 TL (0,9đ) & --- & 11 TN + 1 TL (2,7đ) \\
\hline
\textbf{Tổng cộng} & \textbf{16 câu TN (3,2đ)} & \textbf{13 câu TN (2,6đ)} & \textbf{6 TN + 2 TL (2,2đ)} & \textbf{2 câu TL (1,0đ)} & \textbf{35 TN + 4 TL (10đ)} \\
\hline
\textbf{Tỉ lệ phần trăm} & \textbf{32\%} & \textbf{26\%} & \textbf{22\%} & \textbf{10\%} & \textbf{100\%} \\
\hline
\end{tabularx}
\end{center}

\subsection*{3. Bản đặc tả ma trận đề kiểm tra}
\begin{itemize}[leftmargin=1.5em]
    \item \textbf{Chương I (Lượng giác):}
    \begin{itemize}
        \item \textit{Nhận biết:} Đổi đơn vị độ - radian; dấu giá trị lượng giác; công thức cơ bản ($\sin^2 x + \cos^2 x = 1$); công thức cộng, nhân đôi; tập xác định hàm số lượng giác cơ bản; công thức nghiệm phương trình $\sin x = m, \cos x = m$.
        \item \textit{Thông hiểu:} Tính giá trị lượng giác khi biết một giá trị; rút gọn biểu thức; tìm chu kì hàm số; giải phương trình lượng giác cơ bản.
        \item \textit{Vận dụng \& VDC:} Giải phương trình quy về bậc nhất/bậc hai; tìm số nghiệm thuộc khoảng; bài toán thực tế dao động điều hòa.
    \end{itemize}
    \item \textbf{Chương II (Dãy số, Cấp số):}
    \begin{itemize}
        \item \textit{Nhận biết:} Nhận biết dãy số tăng/giảm; nhận biết cấp số cộng (công sai $d$), cấp số nhân (công bội $q$); công thức số hạng tổng quát $u_n$.
        \item \textit{Thông hiểu:} Tìm số hạng thứ $k$; tìm công sai/công bội; kiểm tra 3 số lập thành cấp số; tính tổng $S_n$.
        \item \textit{Vận dụng \& VDC:} Giải hệ phương trình tìm $u_1, d, q$; bài toán thực tế lãi kép ngân hàng, sự phân đôi tế bào.
    \end{itemize}
    \item \textbf{Chương III (Thống kê ghép nhóm):}
    \begin{itemize}
        \item \textit{Nhận biết:} Nhận biết khoảng nửa mở $[a_i; a_{i+1})$; xác định giá trị đại diện $c_i$; tính tần số tương đối.
        \item \textit{Thông hiểu:} Nhận biết nhóm chứa mốt, nhóm chứa trung vị, nhóm chứa tứ phân vị $Q_1, Q_3$.
        \item \textit{Vận dụng:} Tính số trung bình cộng $\overline{x}$, trung vị $M_e$, mốt $M_o$ và khoảng tứ phân vị $\Delta_Q$; so sánh dữ liệu thực tế.
    \end{itemize}
\end{itemize}

\newpage

\section*{III. ĐỀ KIỂM TRA MINH HỌA GIỮA HỌC KÌ I}

\subsection*{PHẦN I: TRẮC NGHIỆM KHÁCH QUAN (35 CÂU --- 7,0 ĐIỂM)}
\textit{Mỗi câu trắc nghiệm học sinh chọn đúng được 0,2 điểm.}

\begin{enumerate}[leftmargin=1.5em]
    \item Cho góc lượng giác có số đo $\alpha = \dfrac{\pi}{3}$. Số đo theo đơn vị độ của góc đó là:
    \begin{center}
    \textbf{A.} $30^\circ$. \qquad \textbf{B.} $45^\circ$. \qquad \textbf{C.} $60^\circ$. \qquad \textbf{D.} $90^\circ$.
    \end{center}

    \item Điểm biểu diễn của góc lượng giác $\alpha = \dfrac{3\pi}{4}$ nằm ở góc phần tư thứ mấy?
    \begin{center}
    \textbf{A.} Thứ I. \qquad \textbf{B.} Thứ II. \qquad \textbf{C.} Thứ III. \qquad \textbf{D.} Thứ IV.
    \end{center}

    \item Khẳng định nào sau đây là \textbf{đúng} với mọi góc $\alpha$?
    \begin{center}
    \textbf{A.} $\sin^2\alpha - \cos^2\alpha = 1$. \qquad \textbf{B.} $\sin^2\alpha + \cos^2\alpha = 1$. \qquad \textbf{C.} $\tan\alpha \cdot \cot\alpha = -1$. \qquad \textbf{D.} $1 + \tan^2\alpha = \dfrac{1}{\sin^2\alpha}$.
    \end{center}

    \item Công thức nhân đôi nào sau đây là \textbf{đúng}?
    \begin{center}
    \textbf{A.} $\sin 2a = 2\sin a \cos a$. \qquad \textbf{B.} $\cos 2a = \cos^2 a + \sin^2 a$. \qquad \textbf{C.} $\sin 2a = \sin a \cos a$. \qquad \textbf{D.} $\cos 2a = 2\cos a - 1$.
    \end{center}

    \item Tập xác định của hàm số $y = \tan x$ là:
    \begin{center}
    \textbf{A.} $D = \mathbb{R} \setminus \{k\pi, k \in \mathbb{Z}\}$. \qquad \textbf{B.} $D = \mathbb{R} \setminus \left\{\dfrac{\pi}{2} + k\pi, k \in \mathbb{Z}\right\}$. \qquad \textbf{C.} $D = \mathbb{R}$. \qquad \textbf{D.} $D = \mathbb{R} \setminus \left\{\dfrac{\pi}{2} + k2\pi, k \in \mathbb{Z}\right\}$.
    \end{center}

    \item Nghiệm của phương trình $\sin x = 0$ là:
    \begin{center}
    \textbf{A.} $x = k2\pi, k \in \mathbb{Z}$. \qquad \textbf{B.} $x = \dfrac{\pi}{2} + k\pi, k \in \mathbb{Z}$. \qquad \textbf{C.} $x = k\pi, k \in \mathbb{Z}$. \qquad \textbf{D.} $x = \pi + k2\pi, k \in \mathbb{Z}$.
    \end{center}

    \item Nghiệm của phương trình $\cos x = \dfrac{1}{2}$ là:
    \begin{center}
    \textbf{A.} $x = \pm \dfrac{\pi}{3} + k2\pi$. \qquad \textbf{B.} $x = \pm \dfrac{\pi}{6} + k2\pi$. \qquad \textbf{C.} $x = \dfrac{\pi}{3} + k\pi$. \qquad \textbf{D.} $x = \pm \dfrac{2\pi}{3} + k2\pi$.
    \end{center}

    \item Chu kì tuần hoàn của hàm số $y = \sin x$ là:
    \begin{center}
    \textbf{A.} $\pi$. \qquad \textbf{B.} $2\pi$. \qquad \textbf{C.} $3\pi$. \qquad \textbf{D.} $4\pi$.
    \end{center}

    \item Cho dãy số $(u_n)$ với $u_n = 2n + 1$. Số hạng thứ $3$ của dãy số là:
    \begin{center}
    \textbf{A.} $u_3 = 5$. \qquad \textbf{B.} $u_3 = 6$. \qquad \textbf{C.} $u_3 = 7$. \qquad \textbf{D.} $u_3 = 8$.
    \end{center}

    \item Trong các dãy số sau, dãy số nào là cấp số cộng?
    \begin{center}
    \textbf{A.} $1, 3, 6, 9, 12$. \qquad \textbf{B.} $2, 4, 8, 16, 32$. \qquad \textbf{C.} $3, 7, 11, 15, 19$. \qquad \textbf{D.} $1, -1, 1, -1, 1$.
    \end{center}

    \item Cho cấp số cộng $(u_n)$ có $u_1 = 3$ và công sai $d = 4$. Số hạng thứ $4$ của cấp số cộng là:
    \begin{center}
    \textbf{A.} $u_4 = 15$. \qquad \textbf{B.} $u_4 = 12$. \qquad \textbf{C.} $u_4 = 16$. \qquad \textbf{D.} $u_4 = 19$.
    \end{center}

    \item Cho cấp số nhân $(u_n)$ có $u_1 = 2$ và công bội $q = 3$. Số hạng thứ $3$ của cấp số nhân là:
    \begin{center}
    \textbf{A.} $u_3 = 18$. \qquad \textbf{B.} $u_3 = 54$. \qquad \textbf{C.} $u_3 = 8$. \qquad \textbf{D.} $u_3 = 12$.
    \end{center}

    \item Công thức tính tổng $n$ số hạng đầu tiên của cấp số cộng là:
    \begin{center}
    \textbf{A.} $S_n = n(u_1 + u_n)$. \qquad \textbf{B.} $S_n = \dfrac{n(u_1 + u_n)}{2}$. \qquad \textbf{C.} $S_n = \dfrac{u_1 + u_n}{2}$. \qquad \textbf{D.} $S_n = \dfrac{n(u_1 + d)}{2}$.
    \end{center}

    \item Công thức tính tổng $n$ số hạng đầu tiên của cấp số nhân với công bội $q \neq 1$ là:
    \begin{center}
    \textbf{A.} $S_n = \dfrac{u_1(1 - q^n)}{1 - q}$. \qquad \textbf{B.} $S_n = \dfrac{u_1(1 - q)}{1 - q^n}$. \qquad \textbf{C.} $S_n = u_1(1 - q^n)$. \qquad \textbf{D.} $S_n = \dfrac{n(u_1 + u_n)}{2}$.
    \end{center}

    \item Cho nhóm số liệu $[150; 160)$. Giá trị đại diện của nhóm này là:
    \begin{center}
    \textbf{A.} $150$. \qquad \textbf{B.} $155$. \qquad \textbf{C.} $160$. \qquad \textbf{D.} $10$.
    \end{center}

    \item Độ dài của nhóm số liệu $[40; 60)$ là:
    \begin{center}
    \textbf{A.} $10$. \qquad \textbf{B.} $20$. \qquad \textbf{C.} $50$. \qquad \textbf{D.} $100$.
    \end{center}

    \item Một mẫu số liệu ghép nhóm có $n = 40$ quan sát. Tần số tích lũy $cf_p$ để xác định nhóm chứa trung vị $M_e$ cần thỏa mãn:
    \begin{center}
    \textbf{A.} $cf_p \ge 10$. \qquad \textbf{B.} $cf_p \ge 20$. \qquad \textbf{C.} $cf_p \ge 30$. \qquad \textbf{D.} $cf_p \ge 40$.
    \end{center}

    \item Nhóm chứa mốt của mẫu số liệu ghép nhóm là nhóm:
    \begin{center}
    \textbf{A.} Có tần số nhỏ nhất. \qquad \textbf{B.} Có độ dài lớn nhất. \qquad \textbf{C.} Có tần số lớn nhất. \qquad \textbf{D.} Ở vị trí chính giữa bảng.
    \end{center}

    \item Cho $\sin\alpha = \dfrac{3}{5}$ với $0 < \alpha < \dfrac{\pi}{2}$. Giá trị của $\cos\alpha$ bằng:
    \begin{center}
    \textbf{A.} $-\dfrac{4}{5}$. \qquad \textbf{B.} $\dfrac{4}{5}$. \qquad \textbf{C.} $\dfrac{16}{25}$. \qquad \textbf{D.} $\pm\dfrac{4}{5}$.
    \end{center}

    \item Rút gọn biểu thức $A = \cos(\pi - x) + \sin\left(\dfrac{\pi}{2} - x\right)$ ta được:
    \begin{center}
    \textbf{A.} $A = 0$. \qquad \textbf{B.} $A = 2\cos x$. \qquad \textbf{C.} $A = -2\cos x$. \qquad \textbf{D.} $A = 2\sin x$.
    \end{center}

    \item Số nghiệm của phương trình $\sin x = \dfrac{1}{2}$ trên đoạn $[0; 2\pi]$ là:
    \begin{center}
    \textbf{A.} $1$. \qquad \textbf{B.} $2$. \qquad \textbf{C.} $3$. \qquad \textbf{D.} $4$.
    \end{center}

    \item Cho cấp số cộng $(u_n)$ có $u_1 = -2$ và $d = 3$. Số $43$ là số hạng thứ mấy của cấp số cộng?
    \begin{center}
    \textbf{A.} Thứ $14$. \qquad \textbf{B.} Thứ $15$. \qquad \textbf{C.} Thứ $16$. \qquad \textbf{D.} Thứ $17$.
    \end{center}

    \item Cho cấp số nhân $(u_n)$ có $u_1 = 3, q = 2$. Tổng của $6$ số hạng đầu tiên ($S_6$) bằng:
    \begin{center}
    \textbf{A.} $93$. \qquad \textbf{B.} $189$. \qquad \textbf{C.} $192$. \qquad \textbf{D.} $381$.
    \end{center}

    \item Ba số $x, 5, 11$ theo thứ tự lập thành một cấp số cộng. Giá trị của $x$ là:
    \begin{center}
    \textbf{A.} $-1$. \qquad \textbf{B.} $1$. \qquad \textbf{C.} $2$. \qquad \textbf{D.} $3$.
    \end{center}

    \item Ba số dương $2, x, 18$ theo thứ tự lập thành một cấp số nhân. Giá trị của $x$ là:
    \begin{center}
    \textbf{A.} $6$. \qquad \textbf{B.} $9$. \qquad \textbf{C.} $10$. \qquad \textbf{D.} $12$.
    \end{center}

    \item Bảng thống kê điểm môn Toán của một lớp: nhóm $[6; 8)$ có $20$ học sinh trên tổng số $40$ học sinh. Tần số tương đối của nhóm này là:
    \begin{center}
    \textbf{A.} $20\%$. \qquad \textbf{B.} $40\%$. \qquad \textbf{C.} $50\%$. \qquad \textbf{D.} $60\%$.
    \end{center}

    \item Cho mẫu số liệu ghép nhóm có $n = 50$, nhóm chứa trung vị là $[160; 165)$, tần số $m = 15$, tần số tích lũy nhóm đứng trước là $cf = 18$. Trung vị $M_e$ xấp xỉ bằng:
    \begin{center}
    \textbf{A.} $162{,}33$. \qquad \textbf{B.} $163{,}50$. \qquad \textbf{C.} $164{,}12$. \qquad \textbf{D.} $161{,}85$.
    \end{center}

    \item Giá trị lớn nhất của hàm số $y = 3\cos x - 2$ là:
    \begin{center}
    \textbf{A.} $1$. \qquad \textbf{B.} $3$. \qquad \textbf{C.} $5$. \qquad \textbf{D.} $-2$.
    \end{center}

    \item Nghiệm của phương trình $\tan x = \sqrt{3}$ là:
    \begin{center}
    \textbf{A.} $x = \dfrac{\pi}{3} + k\pi, k \in \mathbb{Z}$. \qquad \textbf{B.} $x = \dfrac{\pi}{6} + k\pi, k \in \mathbb{Z}$. \qquad \textbf{C.} $x = \dfrac{\pi}{3} + k2\pi, k \in \mathbb{Z}$. \qquad \textbf{D.} $x = \pm\dfrac{\pi}{3} + k\pi, k \in \mathbb{Z}$.
    \end{center}

    \item Cho dãy số $(u_n)$ có $u_n = \dfrac{n - 1}{n + 1}$. Khẳng định nào sau đây là \textbf{đúng}?
    \begin{center}
    \textbf{A.} Dãy số giảm. \qquad \textbf{B.} Dãy số tăng. \qquad \textbf{C.} Dãy số không tăng không giảm. \qquad \textbf{D.} Dãy số không bị chặn.
    \end{center}

    \item Một rạp hát có hàng ghế đầu tiên gồm $15$ ghế, mỗi hàng ghế tiếp theo nhiều hơn hàng ghế ngay trước nó $2$ ghế. Rạp có $20$ hàng ghế. Tổng số ghế của rạp hát là:
    \begin{center}
    \textbf{A.} $680$. \qquad \textbf{B.} $700$. \qquad \textbf{C.} $640$. \qquad \textbf{D.} $720$.
    \end{center}

    \item Khảo sát điểm số của $30$ học sinh cho các nhóm: $[4; 6)$ có $8$ học sinh, $[6; 8)$ có $14$ học sinh, $[8; 10]$ có $8$ học sinh. Số trung bình cộng $\overline{x}$ là:
    \begin{center}
    \textbf{A.} $6{,}8$. \qquad \textbf{B.} $7{,}0$. \qquad \textbf{C.} $7{,}2$. \qquad \textbf{D.} $7{,}4$.
    \end{center}

    \item Cho mẫu số liệu ghép nhóm có $Q_1 = 12$ và $Q_3 = 20$. Khoảng tứ phân vị $\Delta_Q$ bằng:
    \begin{center}
    \textbf{A.} $32$. \qquad \textbf{B.} $16$. \qquad \textbf{C.} $8$. \qquad \textbf{D.} $4$.
    \end{center}

    \item Phương trình $\cos^2 x - 3\cos x + 2 = 0$ có nghiệm là:
    \begin{center}
    \textbf{A.} $x = k2\pi, k \in \mathbb{Z}$. \qquad \textbf{B.} $x = \dfrac{\pi}{2} + k\pi, k \in \mathbb{Z}$. \qquad \textbf{C.} $x = \pi + k2\pi, k \in \mathbb{Z}$. \qquad \textbf{D.} Vô nghiệm.
    \end{center}

    \item Bác An gửi tiết kiệm $50$ triệu đồng vào ngân hàng với lãi suất $6\%$/năm theo hình thức lãi kép. Sau $3$ năm, tổng số tiền bác An nhận được (làm tròn đến nghìn đồng) là:
    \begin{center}
    \textbf{A.} $59\,551\,000$ đồng. \qquad \textbf{B.} $59\,000\,000$ đồng. \qquad \textbf{C.} $58\,500\,000$ đồng. \qquad \textbf{D.} $60\,120\,000$ đồng.
    \end{center}
\end{enumerate}

\subsection*{PHẦN II: TỰ LUẬN (4 BÀI --- 3,0 ĐIỂM)}

\paragraph{Bài 1 (0,75 điểm):}
Giải phương trình lượng giác:
$$\sqrt{3}\tan\left(2x - \dfrac{\pi}{4}\right) - 1 = 0$$

\paragraph{Bài 2 (0,75 điểm):}
Tìm số hạng đầu $u_1$ và công sai $d$ của cấp số cộng $(u_n)$ biết:
$$\begin{cases} u_2 + u_5 - u_3 = 10 \\ u_1 + u_6 = 17 \end{cases}$$

\paragraph{Bài 3 (1,0 điểm):}
Bảng số liệu sau ghi lại thời gian (tính bằng phút) đi từ nhà đến trường của $40$ học sinh:
\begin{center}
\begin{tabular}{|c|c|c|c|c|}
\hline
\textbf{Thời gian (phút)} & $[10; 15)$ & $[15; 20)$ & $[20; 25)$ & $[25; 30]$ \\
\hline
\textbf{Số học sinh ($m_i$)} & $6$ & $16$ & $12$ & $6$ \\
\hline
\end{tabular}
\end{center}
\begin{itemize}
    \item a) Tính số trung bình cộng $\overline{x}$ của mẫu số liệu ghép nhóm trên.
    \item b) Tính trung vị $M_e$ và mốt $M_o$ của mẫu số liệu.
\end{itemize}

\paragraph{Bài 4 (0,5 điểm - Vận dụng cao):}
Cho ba số $a, b, c$ theo thứ tự lập thành một cấp số nhân có tích $a \cdot b \cdot c = 216$. Đồng thời $a, b - 2, c$ theo thứ tự lập thành một cấp số cộng. Tìm ba số $a, b, c$.

\newpage

\section*{IV. ĐÁP ÁN VÀ THANG ĐIỂM CHI TIẾT}

\subsection*{1. Đáp án phần Trắc nghiệm (35 câu)}
\begin{center}
\begin{tabular}{|c|c|c|c|c|c|c|c|c|c|}
\hline
\textbf{1. C} & \textbf{2. B} & \textbf{3. B} & \textbf{4. A} & \textbf{5. B} & \textbf{6. C} & \textbf{7. A} & \textbf{8. B} & \textbf{9. C} & \textbf{10. C} \\
\hline
\textbf{11. A} & \textbf{12. A} & \textbf{13. B} & \textbf{14. A} & \textbf{15. B} & \textbf{16. B} & \textbf{17. B} & \textbf{18. C} & \textbf{19. B} & \textbf{20. A} \\
\hline
\textbf{21. B} & \textbf{22. C} & \textbf{23. B} & \textbf{24. A} & \textbf{25. A} & \textbf{26. C} & \textbf{27. A} & \textbf{28. A} & \textbf{29. A} & \textbf{30. B} \\
\hline
\textbf{31. A} & \textbf{32. B} & \textbf{33. C} & \textbf{34. A} & \textbf{35. A} & & & & & \\
\hline
\end{tabular}
\end{center}

\subsection*{2. Hướng dẫn chấm phần Tự luận}

\begin{center}
\begin{tabularx}{\linewidth}{|p{2.5cm}|X|c|}
\hline
\textbf{Bài} & \textbf{Nội dung đáp án} & \textbf{Điểm} \\
\hline
\textbf{Bài 1} & $\sqrt{3}\tan\left(2x - \dfrac{\pi}{4}\right) - 1 = 0 \iff \tan\left(2x - \dfrac{\pi}{4}\right) = \dfrac{1}{\sqrt{3}}$ & 0,25đ \\
\cline{2-3}
(0,75 điểm) & $\iff \tan\left(2x - \dfrac{\pi}{4}\right) = \tan\left(\dfrac{\pi}{6}\right) \iff 2x - \dfrac{\pi}{4} = \dfrac{\pi}{6} + k\pi$ ($k \in \mathbb{Z}$) & 0,25đ \\
\cline{2-3}
& $\iff 2x = \dfrac{5\pi}{12} + k\pi \iff x = \dfrac{5\pi}{24} + k\dfrac{\pi}{2}$ ($k \in \mathbb{Z}$). & 0,25đ \\
\hline
\textbf{Bài 2} & Quy đổi về $u_1$ và $d$: $\begin{cases} (u_1 + d) + (u_1 + 4d) - (u_1 + 2d) = 10 \\ u_1 + (u_1 + 5d) = 17 \end{cases}$ & 0,25đ \\
\cline{2-3}
(0,75 điểm) & $\iff \begin{cases} u_1 + 3d = 10 \\ 2u_1 + 5d = 17 \end{cases}$ & 0,25đ \\
\cline{2-3}
& Giải hệ phương trình được: $u_1 = 1$ và $d = 3$. & 0,25đ \\
\hline
\textbf{Bài 3a} & Bảng giá trị đại diện: $c_1 = 12{,}5; c_2 = 17{,}5; c_3 = 22{,}5; c_4 = 27{,}5$. & 0,25đ \\
\cline{2-3}
(0,5 điểm) & $\overline{x} = \dfrac{6 \cdot 12{,}5 + 16 \cdot 17{,}5 + 12 \cdot 22{,}5 + 6 \cdot 27{,}5}{40} = \dfrac{75 + 280 + 270 + 165}{40} = \dfrac{790}{40} = 19{,}75$. & 0,25đ \\
\hline
\textbf{Bài 3b} & Cỡ mẫu $n = 40 \implies n/2 = 20$. Tần số tích lũy: $6, 22, 34, 40 \implies$ Nhóm chứa trung vị là $[15; 20)$. & 0,25đ \\
\cline{2-3}
(0,5 điểm) & $M_e = 15 + \dfrac{20 - 6}{16} \cdot 5 = 15 + \dfrac{70}{16} = 19{,}375$ (phút).\\
& Nhóm chứa mốt là $[15; 20)$ ($m_2 = 16$).\\
& $M_o = 15 + \dfrac{16 - 6}{(16 - 6) + (16 - 12)} \cdot 5 = 15 + \dfrac{10}{14} \cdot 5 \approx 18{,}57$ (phút). & 0,25đ \\
\hline
\textbf{Bài 4} & $a, b, c$ là CSN $\implies a \cdot c = b^2$. Tích $a \cdot b \cdot c = b^3 = 216 \implies b = 6$. & 0,25đ \\
\cline{2-3}
(0,5 điểm) & Ba số $a, 4, c$ lập thành CSC $\implies a + c = 2 \cdot 4 = 8 \implies c = 8 - a$.\\
& Ta có: $a(8 - a) = 36 \iff a^2 - 8a + 36 = 0$ (phương trình có $\Delta' = 16 - 36 = -20 < 0$, vô nghiệm).\\
& \textit{(Nếu đề bài $b + 2$: $a+c = 16 \implies a(16 - a) = 36 \implies a = 2$ hoặc $a = 18$)}. & 0,25đ \\
\hline
\end{tabularx}
\end{center}

\newpage

\section*{V. TIẾN TRÌNH TỔ CHỨC KIỂM TRA ĐÁNH GIÁ (90 PHÚT)}

\begin{center}
    \framebox[1.0\textwidth]{\parbox{0.96\textwidth}{
        \textbf{CÁC BƯỚC TỔ CHỨC KIỂM TRA ĐÁNH GIÁ TẠI LỚP:}\\[0.3em]
        $\bullet$ \textbf{Bước 1 (05 phút): Ổn định và chuẩn bị:}\\
        - Giám thị gọi tên học sinh theo danh sách số báo danh, bố trí ngồi xen kẽ.\\
        - Nhắc nhở quy chế thi cử: Toàn bộ cặp sách, tài liệu, điện thoại thông minh để phía trước lớp. Học sinh chỉ được mang bút viết, thước kẻ và máy tính cầm tay được phép theo quy định của Bộ GD\&ĐT.\\[0.3em]
        $\bullet$ \textbf{Bước 2 (05 phút): Phát đề kiểm tra:}\\
        - Giám thị phát đề thi trắc nghiệm kết hợp tự luận đến từng học sinh.\\
        - Yêu cầu học sinh kiểm tra số trang đề thi, điền đầy đủ Họ tên, Lớp, Số báo danh, Mã đề thi vào giấy làm bài.\\[0.3em]
        $\bullet$ \textbf{Bước 3 (75 phút): Học sinh làm bài thi nghiêm túc:}\\
        - Giám thị bao quát phòng thi, duy trì trật tự và tính nghiêm túc tuyệt đối.\\
        - Nhắc nhở thời gian làm bài ở các mốc: còn 45 phút, còn 15 phút, còn 5 phút để học sinh hoàn thiện tô phiếu trắc nghiệm và soát bài tự luận.\\[0.3em]
        $\bullet$ \textbf{Bước 4 (05 phút): Thu bài và kết thúc:}\\
        - Hết giờ làm bài, học sinh ngừng bút, nộp bài thi theo thứ tự danh sách.\\
        - Giám thị đếm đủ số bài thi, đối chiếu chữ kí thí sinh, niêm phong bài thi bàn giao về Tổ chuyên môn chấm thi.
    }}
\end{center}

\end{document}
